\documentclass[12pt]{article}
\usepackage[utf8]{inputenc}
\usepackage[T1]{fontenc}
\usepackage{hyperref}
\title{Phuma Language Resources for K-12 Curriculum Development}
\author{Compiled by Claude Sonnet 5.5 (claude-sonnet-5-5) for Andr\'as Kornai}
\date{}
\begin{document}
\maketitle

\section{Introduction}
Phuma (ISO 639-3 ypm; Glottocode phum1235) is a Loloish (Sino-Tibetan) language of southern Yunnan, China, mainly Jinping County near the Vietnam border. Other names found in search results are Black Muji, Hei Muji, Bokha and Phula. Joshua Project lists about 8,700 speakers. Wikitongues classifies the language as threatened (EGIDS 6b). This differs from the ``stable'' status in the reference description. Searches found no standard written form and no Bible translation.

\section{Primary Educational Resources}
No primers, textbooks or school programmes for Phuma turned up in the searches. None are known to exist.

\section{Linguistic and Reference Materials}
\begin{itemize}
\item Wikipedia, Muji language. It treats Phuma as a variety closely related to Muji and notes that the variety names are used inconsistently: \url{https://en.wikipedia.org/wiki/Muji_language}
\item Glottolog entry for Phuma, placed under Bokha-Phuma: \url{https://glottolog.org/resource/languoid/id/phum1235.bigmap.html}
\item Grambank language page (found in results; I did not confirm that it covers Phuma specifically): \url{https://grambank.clld.org/languages/sout2729}
\item Endangered Languages Project bibliography page: \url{https://endangeredlanguages.com/lang/1125/bibliography}
\end{itemize}

\section{Religious and Community Resources}
Joshua Project reports no Bible translation. Its pages are written from a missionary perspective:
\begin{itemize}
\item \url{https://joshuaproject.net/people_groups/20610/CH}
\item \url{https://joshuaproject.net/languages/ypm}
\end{itemize}
No community organisations were found.

\section{Digital and Audio Resources}
Wikitongues has a page for ypm. I could not tell whether it holds any recordings: \url{https://wikitongues.org/languages/ypm/}

\section{Curriculum Development Context}
There is almost no published educational material, and no written standard was found. A curriculum would have to start with documentation, orthography design and recording work together with the community. Most further information would have to come from fieldwork and from specialist literature on Muji and the Southeastern Loloish languages.
\end{document}